\documentclass[10pt,a4paper]{amsart}
\usepackage{amsmath,amssymb,amsthm,graphicx}
\usepackage{verbatim}
\newenvironment{reference}{\comment}{\endcomment}
%\newenvironment{reference}{}{}
\newenvironment{slogan}{\comment}{\endcomment}
\newenvironment{history}{\comment}{\endcomment}

% For commutative diagrams we use Xy-pic
\usepackage[all]{xy}

% We use 2cell for 2-commutative diagrams.
\xyoption{2cell}
% Theorem environments.
%
\theoremstyle{plain}
\newtheorem{theorem}[subsection]{Theorem}
\newtheorem{proposition}[subsection]{Proposition}
\newtheorem{lemma}[subsection]{Lemma}

\theoremstyle{definition}
\newtheorem{definition}[subsection]{Definition}
\newtheorem{example}[subsection]{Example}
\newtheorem{exercise}[subsection]{Exercise}
\newtheorem{situation}[subsection]{Situation}

\theoremstyle{remark}
\newtheorem{remark}[subsection]{Remark}
\newtheorem{remarks}[subsection]{Remarks}

\numberwithin{equation}{subsection}

% Macros
%
\def\lim{\mathop{\mathrm{lim}}\nolimits}
\def\colim{\mathop{\mathrm{colim}}\nolimits}
\def\Spec{\mathop{\mathrm{Spec}}}
\def\Hom{\mathop{\mathrm{Hom}}\nolimits}
\def\Ext{\mathop{\mathrm{Ext}}\nolimits}
\def\SheafHom{\mathop{\mathcal{H}\!\mathit{om}}\nolimits}
\def\SheafExt{\mathop{\mathcal{E}\!\mathit{xt}}\nolimits}
\def\Sch{\mathit{Sch}}
\def\Mor{\mathop{\mathrm{Mor}}\nolimits}
\def\Ob{\mathop{\mathrm{Ob}}\nolimits}
\def\Sh{\mathop{\mathit{Sh}}\nolimits}
\def\NL{\mathop{N\!L}\nolimits}
\def\CH{\mathop{\mathrm{CH}}\nolimits}
\def\proetale{{pro\text{-}\acute{e}tale}}
\def\etale{{\acute{e}tale}}
\def\QCoh{\mathit{QCoh}}
\def\Ker{\mathop{\mathrm{Ker}}}
\def\Im{\mathop{\mathrm{Im}}}
\def\Coker{\mathop{\mathrm{Coker}}}
\def\Coim{\mathop{\mathrm{Coim}}}

% Boxtimes
%
\DeclareMathSymbol{\boxtimes}{\mathbin}{AMSa}{"02}

%
% Macros for moduli stacks/spaces
%
\def\QCohstack{\mathcal{QC}\!\mathit{oh}}
\def\Cohstack{\mathcal{C}\!\mathit{oh}}
\def\Spacesstack{\mathcal{S}\!\mathit{paces}}
\def\Quotfunctor{\mathrm{Quot}}
\def\Hilbfunctor{\mathrm{Hilb}}
\def\Curvesstack{\mathcal{C}\!\mathit{urves}}
\def\Polarizedstack{\mathcal{P}\!\mathit{olarized}}
\def\Complexesstack{\mathcal{C}\!\mathit{omplexes}}
% \Pic is the operator that assigns to X its picard group, usage \Pic(X)
% \Picardstack_{X/B} denotes the Picard stack of X over B
% \Picardfunctor_{X/B} denotes the Picard functor of X over B
\def\Pic{\mathop{\mathrm{Pic}}\nolimits}
\def\Picardstack{\mathcal{P}\!\mathit{ic}}
\def\Picardfunctor{\mathrm{Pic}}
\def\Deformationcategory{\mathcal{D}\!\mathit{ef}}
\usepackage{cleveref}
\crefname{originalreference}{original reference}{original references}
\Crefname{originalreference}{Original reference}{Original references}
\makeatletter
\newlabel{example-glue-points}{{0C1J}{source}{}{sourcepointer-1}{}}
\newlabel{example-glue-points@cref}{{[originalreference][0][]0C1J}{[1][0][]source}{}{}{}}
\makeatother
\crefname{originalreference}{original reference}{original references}
\Crefname{originalreference}{Original reference}{Original references}
\makeatletter
\newlabel{example-squish-tangent-vector}{{0C1K}{source}{}{sourcepointer-2}{}}
\newlabel{example-squish-tangent-vector@cref}{{[originalreference][0][]0C1K}{[1][0][]source}{}{}{}}
\makeatother
\crefname{originalreference}{original reference}{original references}
\Crefname{originalreference}{Original reference}{Original references}
\makeatletter
\newlabel{lemma-factor-almost-isomorphism}{{0C1L}{source}{}{sourcepointer-3}{}}
\newlabel{lemma-factor-almost-isomorphism@cref}{{[originalreference][0][]0C1L}{[1][0][]source}{}{}{}}
\makeatother
\crefname{originalreference}{original reference}{original references}
\Crefname{originalreference}{Original reference}{Original references}
\makeatletter
\newlabel{lemma-glue-points}{{0C1M}{source}{}{sourcepointer-4}{}}
\newlabel{lemma-glue-points@cref}{{[originalreference][0][]0C1M}{[1][0][]source}{}{}{}}
\makeatother
\crefname{originalreference}{original reference}{original references}
\Crefname{originalreference}{Original reference}{Original references}
\makeatletter
\newlabel{lemma-multicross}{{0C1X}{source}{}{sourcepointer-5}{}}
\newlabel{lemma-multicross@cref}{{[originalreference][0][]0C1X}{[1][0][]source}{}{}{}}
\makeatother
\crefname{originalreference}{original reference}{original references}
\Crefname{originalreference}{Original reference}{Original references}
\makeatletter
\newlabel{lemma-squish-tangent-vector}{{0C1N}{source}{}{sourcepointer-6}{}}
\newlabel{lemma-squish-tangent-vector@cref}{{[originalreference][0][]0C1N}{[1][0][]source}{}{}{}}
\makeatother
\crefname{originalreference}{original reference}{original references}
\Crefname{originalreference}{Original reference}{Original references}
\makeatletter
\newlabel{lemma-torsion-picard-smooth-projective}{{0C1Z}{source}{}{sourcepointer-7}{}}
\newlabel{lemma-torsion-picard-smooth-projective@cref}{{[originalreference][0][]0C1Z}{[1][0][]source}{}{}{}}
\makeatother
\crefname{originallemma}{original lemma}{original lemmas}
\Crefname{originallemma}{Original lemma}{Original lemmas}
\makeatletter
\newlabel{varieties-lemma-prepare-delta-invariant}{{0C1R}{source}{}{sourcepointer-8}{}}
\newlabel{varieties-lemma-prepare-delta-invariant@cref}{{[originallemma][0][]0C1R}{[1][0][]source}{}{}{}}
\makeatother
\title{l-torsion of the Picard group of a reduced connected proper curve is bounded by genus plus geometric genus, with equality exactly for multicross singularities}
\author{The Stacks Project Authors}
\date{stacks-project:curves.tex:local-snapshot-20261012; GNU FDL 1.2 or later}
\begin{document}
\maketitle
\noindent\textit{Editorial note.} All delivered mathematical text is the author's own, copied from the pinned Stacks Project chapter file curves.tex; no statement or proof step was rewritten, reordered or removed.\par\medskip
\noindent\textit{Editorial note.} The counted proof is the author's complete argument: the three-step factorization with its Picard sequences and cohomology changes, the l-divisibility computation, the endpoint values on the normalization, the counting bound, and the two directions of the equality characterization are retained verbatim, including the author's remark that the final gluing of local sequences is given as a hint with some details omitted.\par\medskip
\noindent\textit{Editorial note.} The label proposition-torsion-picard-reduced-proper names the counted statement, and the proof environment immediately following it in the source is delivered verbatim as its proof.\par\medskip
\noindent\textit{Editorial note.} The counted statement is carried by the author's own proposition environment in the chapter source.\par\medskip
\noindent\textit{Allowed standard prerequisites (editorial).}\begin{itemize}
\item Factorization of the normalization into elementary modifications. Same chapter, Lemma lemma-factor-almost-isomorphism with Lemmas lemma-glue-points and lemma-squish-tangent-vector, and Examples example-glue-points, example-squish-tangent-vector. Supplies the sequence X\_n -> ... -> X\_0 = X where each step glues two points or squishes a tangent vector, with the induced exact sequences on Picard and the changes of h\textasciicircum{}0, h\textasciicircum{}1.
\item Torsion Picard of a smooth projective curve. Same chapter, Lemma lemma-torsion-picard-smooth-projective; Varieties, Lemma varieties-lemma-prepare-delta-invariant. Gives the endpoint values on the normalization X\_n = disjoint union of smooth projective curves: h\textasciicircum{}1\_n = g\_geom and dim Pic(X\_n)[l] = 2 g\_geom.
\end{itemize}\medskip
\ifdefined\hypertarget\hypertarget{sourcepointer-1}{}\fi
\noindent\textit{Source reference.} Original reference 0C1J: sources/stacks-curves/pinned-original-source.tex, source lines 3202--3202. Prerequisite pointer invoked at proof line(s) 3749; it names a standard fact of the same chapter that the counted argument uses, and every step of the counted proof built on it is retained verbatim in the delivered body.\par\medskip
\ifdefined\hypertarget\hypertarget{sourcepointer-2}{}\fi
\noindent\textit{Source reference.} Original reference 0C1K: sources/stacks-curves/pinned-original-source.tex, source lines 3228--3228. Prerequisite pointer invoked at proof line(s) 3749; it names a standard fact of the same chapter that the counted argument uses, and every step of the counted proof built on it is retained verbatim in the delivered body.\par\medskip
\ifdefined\hypertarget\hypertarget{sourcepointer-3}{}\fi
\noindent\textit{Source reference.} Original reference 0C1L: sources/stacks-curves/pinned-original-source.tex, source lines 3257--3257. Prerequisite pointer invoked at proof line(s) 3722; it names a standard fact of the same chapter that the counted argument uses, and every step of the counted proof built on it is retained verbatim in the delivered body.\par\medskip
\ifdefined\hypertarget\hypertarget{sourcepointer-4}{}\fi
\noindent\textit{Source reference.} Original reference 0C1M: sources/stacks-curves/pinned-original-source.tex, source lines 3323--3323. Prerequisite pointer invoked at proof line(s) 3725; it names a standard fact of the same chapter that the counted argument uses, and every step of the counted proof built on it is retained verbatim in the delivered body.\par\medskip
\ifdefined\hypertarget\hypertarget{sourcepointer-5}{}\fi
\noindent\textit{Source reference.} Original reference 0C1X: sources/stacks-curves/pinned-original-source.tex, source lines 3478--3478. Prerequisite pointer invoked at proof line(s) 3776; it names a standard fact of the same chapter that the counted argument uses, and every step of the counted proof built on it is retained verbatim in the delivered body.\par\medskip
\ifdefined\hypertarget\hypertarget{sourcepointer-6}{}\fi
\noindent\textit{Source reference.} Original reference 0C1N: sources/stacks-curves/pinned-original-source.tex, source lines 3364--3364. Prerequisite pointer invoked at proof line(s) 3725; it names a standard fact of the same chapter that the counted argument uses, and every step of the counted proof built on it is retained verbatim in the delivered body.\par\medskip
\ifdefined\hypertarget\hypertarget{sourcepointer-7}{}\fi
\noindent\textit{Source reference.} Original reference 0C1Z: sources/stacks-curves/pinned-original-source.tex, source lines 3611--3611. Prerequisite pointer invoked at proof line(s) 3766; it names a standard fact of the same chapter that the counted argument uses, and every step of the counted proof built on it is retained verbatim in the delivered body.\par\medskip
\ifdefined\hypertarget\hypertarget{sourcepointer-8}{}\fi
\noindent\textit{Source reference.} Original lemma 0C1R: sources/stacks-chow/tags.txt, source lines 15131--15131. Prerequisite pointer invoked at proof line(s) 3717; it names a standard fact of a sibling chapter that the counted argument uses, and every step of the counted proof built on it is retained verbatim in the delivered body.\par\medskip
\begin{proposition}
\label{proposition-torsion-picard-reduced-proper}
Let $k$ be an algebraically closed field. Let $X$ be a proper scheme over $k$
which is reduced, connected, and has dimension $1$. Let $g$ be the genus
of $X$ and let $g_{geom}$ be the sum of the geometric genera of the
irreducible components of $X$. For any prime $\ell$ different from
the characteristic of $k$ we have
$$
\dim_{\mathbf{F}_\ell} \Pic(X)[\ell]
\leq g + g_{geom}
$$
and equality holds if and only if all the singularities of $X$
are multicross.
\end{proposition}

\begin{proof}
Let $\nu : X^\nu \to X$ be the normalization
(Varieties, Lemma \ref{varieties-lemma-prepare-delta-invariant}).
Choose a factorization
$$
X^\nu = X_n \to X_{n - 1} \to \ldots \to X_1 \to X_0 = X
$$
as in Lemma \ref{lemma-factor-almost-isomorphism}.
Let us denote $h^0_i = \dim_k H^0(X_i, \mathcal{O}_{X_i})$
and $h^1_i = \dim_k H^1(X_i, \mathcal{O}_{X_i})$.
By Lemmas \ref{lemma-glue-points} and \ref{lemma-squish-tangent-vector}
for each $n > i \geq 0$ we have
one of the following there possibilities
\begin{enumerate}
\item $X_i$ is obtained by glueing $a, b \in X_{i + 1}$
which are on different connected components: in this case
$\Pic(X_i) = \Pic(X_{i + 1})$,
$h^0_{i + 1} = h^0_i + 1$, $h^1_{i + 1} = h^1_i$,
\item $X_i$ is obtained by glueing $a, b \in X_{i + 1}$
which are on the same connected component: in this case
there is a short exact sequence
$$
0 \to k^* \to \Pic(X_i) \to \Pic(X_{i + 1}) \to 0,
$$
and $h^0_{i + 1} = h^0_i$, $h^1_{i + 1} = h^1_i - 1$,
\item $X_i$ is obtained by squishing a tangent vector in $X_{i + 1}$:
in this case there is a short exact sequence
$$
0 \to (k, +) \to \Pic(X_i) \to \Pic(X_{i + 1}) \to 0,
$$
and $h^0_{i + 1} = h^0_i$, $h^1_{i + 1} = h^1_i - 1$.
\end{enumerate}
To prove the statements on dimensions of cohomology groups of the
structure sheaf, use the exact sequences in
Examples \ref{example-glue-points} and \ref{example-squish-tangent-vector}.
Since $k$ is algebraically closed of characteristic prime to $\ell$
we see that $(k, +)$ and $k^*$ are $\ell$-divisible and with
$\ell$-torsion $(k, +)[\ell] = 0$ and $k^*[\ell] \cong \mathbf{F}_\ell$.
Hence
$$
\dim_{\mathbf{F}_\ell} \Pic(X_{i + 1})[\ell] -
\dim_{\mathbf{F}_\ell}\Pic(X_i)[\ell]
$$
is zero, except in case (2) where it is equal to $-1$.
At the end of this process we get the normalization
$X^\nu = X_n$ which is a disjoint union of smooth projective
curves over $k$. Hence we have
\begin{enumerate}
\item $h^1_n = g_{geom}$ and
\item $\dim_{\mathbf{F}_\ell} \Pic(X_n)[\ell] = 2g_{geom}$.
\end{enumerate}
The last equality by Lemma \ref{lemma-torsion-picard-smooth-projective}.
Since $g = h^1_0$ we see that the number of steps of type
(2) and (3) is at most $h^1_0 - h^1_n = g - g_{geom}$.
By our comptation of the differences in ranks we conclude that
$$
\dim_{\mathbf{F}_\ell} \Pic(X)[\ell] \leq
g - g_{geom} + 2g_{geom} = g + g_{geom}
$$
and equality holds if and only if no steps of type (3) occur.
This indeed means that all singularities of $X$ are multicross
by Lemma \ref{lemma-multicross}. Conversely, if all the singularities
are multicross, then Lemma \ref{lemma-multicross} guarantees that
we can find a sequence $X^\nu = X_n \to \ldots \to X_0 = X$
as above such that no steps of type (3) occur in the sequence
and we find equality holds in the lemma (just glue the local sequences
for each point to find one that works for all singular points of $x$;
some details omitted).
\end{proof}

\end{document}
